% from ega2-1-fr.tex
\oldpage[II]{52}

  \item[\rm(ii)] Soit $\mathcal{S}'$ la $\mathcal{O}_Y$-Algèbre graduée
    somme directe de $\mathcal{O}_Y$ et des $\mathcal{S}_n$ ($n\geq0$)~;
    alors $\operatorname{Proj}(\mathcal{S}')$ et
    $\operatorname{Proj}(\mathcal{S})$ sont canoniquement $Y$-isomorphes.
  \item[\rm(iii)] Soit $\mathcal{L}$ un $\mathcal{O}_Y$-Module inversible
    (\hyperref[0.5.4.1-fr]{0, 5.4.1}) et soit $\mathcal{S}_{(\mathcal{L})}$ la
    $\mathcal{O}_Y$-Algèbre graduée somme directe des
    $\mathcal{S}_d\otimes\mathcal{L}^{\otimes d}$ ($d\geq0$)~; alors
    $\operatorname{Proj}(\mathcal{S})$ et
    $\operatorname{Proj}(\mathcal{S}_{(\mathcal{L})})$ sont canoniquement
    $Y$-isomorphes.
\end{enumerate}
\end{proposition}

\begin{proof}
Dans chacun des trois cas, il suffit de définir l'isomorphisme localement sur
$Y$, la vérification des compatibilités avec les opérations de restriction
d'un ouvert à un ouvert plus petit étant triviale. On peut donc supposer $Y$
affine, et alors (i) résulte de (\hyperref[II.2.4.7-fr]{2.4.7, (i)}) et (ii)
de (\hyperref[II.2.4.8-fr]{2.4.8}). En ce qui concerne (iii), lorsqu'on
suppose en outre $\mathcal{L}$ isomorphe à $\mathcal{O}_Y$ (ce qui est permis,
la question étant locale sur $Y$), l'isomorphie de
$\operatorname{Proj}(\mathcal{S})$ et
$\operatorname{Proj}(\mathcal{S}_{(\mathcal{L})})$ est évidente~; pour
définir un isomorphisme \emph{canonique}, soient
$Y=\operatorname{Spec}(A)$, $\mathcal{S}=\widetilde{S}$, où $S$ est une
$A$-algèbre graduée, et soit $c$ un générateur du $A$-module libre $L$ tel
que $\mathcal{L}=\widetilde{L}$~; alors pour tout $n>0$,
$x_n\mapsto x_n\otimes c^{\otimes n}$ est un $A$-isomorphisme de $S_n$ sur
$S_n\otimes L^{\otimes n}$, et ces $A$-isomorphismes définissent un
$A$-isomorphisme d'algèbres graduées
\[
  \varphi_c:S\longrightarrow S_{(L)}
  =\bigoplus_{n\geq0}S_n\otimes L^{\otimes n}.
\]
Soit alors $f\in S_+$, homogène de degré $d$~; pour tout $x\in S_{nd}$, on a
\[
  \frac{x\otimes c^{\otimes nd}}{(f\otimes c^{\otimes d})^n}
  =\frac{x\otimes(\varepsilon c)^{\otimes nd}}
  {(f\otimes(\varepsilon c)^{\otimes d})^n}
\]
pour tout élément inversible $\varepsilon\in A$, ce qui montre que
l'isomorphisme $S_{(f)}\to(S_{(L)})_{(f\otimes c^d)}$ déduit de $\varphi_c$
est \emph{indépendant} du générateur $c$ de $L$ considéré et achève la
démonstration.
\end{proof}

\begin{env}[3.1.9]
\label{II.3.1.9-fr}
Rappelons ((\hyperref[0.4.1.3-fr]{0, 4.1.3}) et
(\hyperref[I.1.3.14-fr]{I, 1.3.14})) que pour que la
$\mathcal{O}_Y$-Algèbre graduée quasi-cohérente $\mathcal{S}$ soit
\emph{engendrée par le $\mathcal{O}_Y$-Module $\mathcal{S}_1$}, il faut et il
suffit qu'il existe un recouvrement $(U_\alpha)$ de $Y$ par des ouverts
affines tels que l'algèbre graduée $\Gamma(U_\alpha,\mathcal{S})$ sur
$\Gamma(U_\alpha,\mathcal{S}_0)$ soit engendrée par l'ensemble
$\Gamma(U_\alpha,\mathcal{S}_1)$ de ses éléments homogènes de degré~$1$. Pour
tout ouvert $V$ de $Y$, $\mathcal{S}|V$ est alors engendrée par le
$(\mathcal{O}_Y|V)$-Module $\mathcal{S}_1|V$.
\end{env}

\begin{proposition}[3.1.10]
\label{II.3.1.10-fr}
Supposons qu'il existe un recouvrement ouvert affine fini $(U_i)$ de $Y$ tel
que chacune des algèbres graduées $\Gamma(U_i,\mathcal{S})$ soit de type fini
sur $\Gamma(U_i,\mathcal{O}_Y)$. Alors il existe $d>0$ tel que
$\mathcal{S}^{(d)}$ soit engendrée par $\mathcal{S}_d$, $\mathcal{S}_d$ étant
un $\mathcal{O}_Y$-Module de type fini.
\end{proposition}

\begin{proof}
En effet, il résulte de (\hyperref[II.2.1.6-fr]{2.1.6, (v)}) que pour chaque
$i$, il existe un entier $m_i$ tel que
$\Gamma(U_i,\mathcal{S}_{nm_i})=(\Gamma(U_i,\mathcal{S}_{m_i}))^n$ pour tout
$n>0$~; il suffit de prendre pour $d$ un multiple commun des $m_i$, compte
tenu de (\hyperref[II.2.1.6-fr]{2.1.6, (i)}).
\end{proof}

\begin{corollary}[3.1.11]
\label{II.3.1.11-fr}
Sous les hypothèses de (\hyperref[II.3.1.10-fr]{3.1.10}),
$\operatorname{Proj}(\mathcal{S})$ est $Y$-isomorphe à un spectre homogène
$\operatorname{Proj}(\mathcal{S}')$, où $\mathcal{S}'$ est une
$\mathcal{O}_Y$-Algèbre graduée engendrée par $\mathcal{S}'_1$,
$\mathcal{S}'_1$ étant supposé être un $\mathcal{O}_Y$-Module de type fini.
\end{corollary}

\begin{proof}
Il suffit en effet de prendre $\mathcal{S}'=\mathcal{S}^{(d)}$, $d$ étant
déterminé par la propriété de (\hyperref[II.3.1.10-fr]{3.1.10}), et
d'appliquer (\hyperref[II.3.1.8-fr]{3.1.8, (i)}).
\end{proof}

\begin{env}[3.1.12]
\label{II.3.1.12-fr}
Si $\mathcal{S}$ est une $\mathcal{O}_Y$-Algèbre graduée quasi-cohérente à
degrés positifs, on sait (\hyperref[I.5.1.1-fr]{I, 5.1.1}) que son
\emph{nilradical} $\mathcal{N}$ est un $\mathcal{O}_Y$-Module
quasi-cohérent~; nous dirons que $\mathcal{N}_+=\mathcal{N}\cap\mathcal{S}_+$
est le \emph{nilradical de $\mathcal{S}_+$}~; c'est un
$\mathcal{S}_0$-Module quasi-cohérent gradué, car on se ramène immédiatement
au cas où $Y$ est affine, et la proposition résulte alors de
(\hyperref[II.2.1.10-fr]{2.1.10}). Pour tout $y\in Y$,
$(\mathcal{N}_+)_y$ est alors le nilradical de
$(\mathcal{S}_+)_y=(\mathcal{S}_y)_+$
(\hyperref[I.5.1.1-fr]{I, 5.1.1}). On dit que la
$\mathcal{O}_Y$-Algèbre graduée $\mathcal{S}$ est \emph{essentiellement
réduite} si $\mathcal{N}_+=0$, ce qui équivaut

